\documentclass[11pt]{article}

% ----- Packages -----
\usepackage[margin=1in]{geometry}   % Page margins
\usepackage{array}                   % Better column control
\usepackage{tabularx}                % Tables with flexible widths
\usepackage{ragged2e}                % Better text wrapping in cells
\usepackage{xcolor}                  % Color support
\usepackage{colortbl}                % Colored table cells/rows
\usepackage{setspace}                % Line spacing
\usepackage{amssymb}                 % Provides \square for checkboxes
\usepackage{enumitem}                % Better list control

% ----- Custom column types for tabularx -----
% L = left-aligned wrappable column
% C = centered wrappable column
\newcolumntype{L}[1]{>{\RaggedRight\arraybackslash}p{#1}}
\newcolumntype{C}[1]{>{\Centering\arraybackslash}p{#1}}
\newcolumntype{Y}{>{\RaggedRight\arraybackslash}X}

% ----- Header row styling -----
\definecolor{headergray}{RGB}{220,220,220}
\newcommand{\headerrow}[1]{\cellcolor{headergray}\textbf{#1}}

\begin{document}

% ==================== TITLE & DESCRIPTION ====================
\begin{center}
	{\LARGE \textbf{Ceramic Mug Evaluation Rubric}} \\[0.5em]
	{\large A Framework For Assessing Brandon's Mugs}
\end{center}

\vspace{0.5em}

\noindent
This rubric is designed to assess what Brandon Blaylock considers to be a good
mug. It's focused on his personal aesthetic and functional considerations. To
use the rubric for an assessment start with a quality (ie. Ergonomics) and move
from Does Not Meet to Fully Meets. The first quality that a piece matches is the
rating used.

\vspace{1.5em}

% ==================== RUBRIC MATRIX ====================
\begin{center}
	\renewcommand{\arraystretch}{1.6}  % Increase row height for readability
	\begin{tabularx}{\textwidth}{| L{2.5cm} | Y | Y | Y |}
		\hline
		\headerrow{Quality}                                                      &
		\headerrow{\Centering Does Not Meet}                                     &
		\headerrow{\Centering Partially Meets}                                   &
		\headerrow{\Centering Fully Meets}                                         \\
		\hline

		% ----- Row 1: Ergonomics -----
		\textbf{Ergonomics}                                                      &
		Uncomfortable handle, sharp edges, chunky or awkward rim.                &
		Hard to wash, holds less than 300ml, unbalanced weigh, wobbles on tablet &
		No obvious awkward physical traits                                         \\
		\hline

		% ----- Row 2: Texture -----
		\textbf{Texture}                                                         &
		No textures or texture seems accidental                                  &
		Some texture but seems inconsistently applied                            &
		Texture is consistent and interesting                                      \\
		\hline

		% ----- Row 3: Glaze Quality -----
		\textbf{Glaze Quality}                                                   &
		Glaze is uneven, pinholed, or crawling.                                  &
		Glaze is consistent but lacks depth or color development                 &
		Glaze is fully developed and multiple glazes have crisp delineations       \\
		\hline

		% ----- Add more rows here following the same pattern -----
		% \textbf{New Quality} &
		%     Description for "does not meet." &
		%     Description for "partially meets." &
		%     Description for "fully meets." \\
		% \hline
	\end{tabularx}
\end{center}

\vspace{1.5em}

% ==================== EVALUATION FORMS ====================

% ----- Checkbox shortcut -----
\newcommand{\cbox}{$\square$}

% ----- A single evaluation form -----
% To add/remove a quality, edit the list inside this command.
% Make sure the qualities here match the rows in the rubric above.
\newcommand{\evalform}{%
	\textbf{Piece Name:} \underline{\hspace{0.55\linewidth}} \\[0.6em]
	\textbf{Evaluation:} \\[0.2em]
	\begin{tabular}{@{}l l@{}}
		Ergonomics:    & \cbox\ None \quad \cbox\ Partial \quad \cbox\ Full \\
		Texture:       & \cbox\ None \quad \cbox\ Partial \quad \cbox\ Full \\
		Glaze Quality: & \cbox\ None \quad \cbox\ Partial \quad \cbox\ Full \\
		% Add more qualities here, e.g.:
		% New Quality:  & \cbox\ None \quad \cbox\ Partial \quad \cbox\ Full \\
	\end{tabular} \\[0.6em]
	\textbf{Notes:} \\[0.2em]
	\rule{\linewidth}{0.4pt} \\[1em]
	\rule{\linewidth}{0.4pt} \\[1em]
	\rule{\linewidth}{0.4pt} \\[1em]
	\rule{\linewidth}{0.4pt}
}

% ----- 2-column x 3-row grid of forms -----
% Each row uses two minipages side-by-side.
% \formgap controls spacing between rows; \colgap controls the gutter between columns.
\newlength{\colgap}\setlength{\colgap}{0.04\textwidth}
\newlength{\formwidth}\setlength{\formwidth}{0.48\textwidth}
\newcommand{\formrow}{%
	\noindent
	\begin{minipage}[t]{\formwidth}\evalform\end{minipage}%
	\hspace{\colgap}%
	\begin{minipage}[t]{\formwidth}\evalform\end{minipage}%
	\par\vspace{3em}
}

\formrow
\formrow
\formrow
\formrow

\end{document}

